% Part: sets-functions-relations
% Chapter: sets
% Section: proofs-about-sets

\documentclass[../../../include/open-logic-section]{subfiles}

\begin{document}

\olfileid{sfr}{set}{prf}
\olsection{సమితుల గురించి నిరూపణలు}

\begin{editorial}
ఈ విభాగం స్థానంలో ఇప్పుడు \verb|content/method/proofs| ఉంది; అందువల్ల అప్రమేయంగా ఈ అధ్యాయంలో దీన్ని చేర్చరు.
\end{editorial}

\begin{explain}
ఇప్పటివరకు పరిచయం చేసిన సమితులు, సంకేతాలు సమితులను నిర్దేశించడానికి, వాటి మధ్య సంబంధాలను వ్యక్తపరచడానికి సౌకర్యమైన సంక్షిప్త రూపాలు ఇస్తాయి. అటువంటి సంబంధాల గురించి వాదనలను నిరూపించాల్సిన అవసరం కూడా తరచూ ఉంటుంది. గణిత నిరూపణలు మీకు పరిచయం లేకపోతే, ఇది కొత్తగా అనిపించవచ్చు. అందుకే సరళ ఉదాహరణను దశలవారీగా చూద్దాం. ఏ సమితులు $X$, $Y$లకైనా $X \cap (X \cup Y) = X$ అని నిరూపిద్దాం. సమితుల మధ్య ఇటువంటి సమానత్వాన్ని ఎలా నిరూపిస్తాం? సమితులను వాటి \tetoken{మూలకాలే}{!!{element}s} పూర్తిగా నిర్ధారిస్తాయని గుర్తుంచుకోండి; అంటే రెండు సమితులకు ఒకే \tetoken{మూలకాలు}{!!{element}s} ఉంటే, ఉన్నప్పుడే అవి సమానమైనవి. కాబట్టి ఇక్కడ (a) $X \cap (X \cup Y)$ యొక్క ప్రతి \tetoken{మూలకం}{!!{element}} కూడా $X$లో ఒక \tetoken{మూలకం}{!!{element}} అని, వ్యతిరేక దిశలో (b) $X$లోని ప్రతి \tetoken{మూలకం}{!!{element}} కూడా $X \cap (X \cup Y)$లో ఒక \tetoken{మూలకం}{!!{element}} అని నిరూపించాలి. మరో మాటలో, (a) $X \cap (X \cup Y) \subseteq X$, (b) $X \subseteq X \cap (X \cup Y)$ రెండింటినీ చూపాలి.

“$X \cap (X \cup Y)$ యొక్క ప్రతి \tetoken{మూలకం}{!!{element}}~$z$, $X$ యొక్క \tetoken{మూలకం}{!!{element}} కూడా” అనే సార్వత్రిక వాదనను నిరూపించడానికి, ముందుగా ఏదైనా $z \in X \cap (X \cup Y)$ ఇచ్చారని అనుకుని, ఆ పరికల్పన నుంచి $z \in X$ అని నిరూపిస్తాం. ఈ నమూనాను “సార్వత్రిక షరతులతో కూడిన నిరూపణ”గా మీరు తెలుసుకుని ఉండవచ్చు. ఈ నిరూపణలో పరికల్పన, నిష్కర్షలలోని నిర్వచనాలను కూడా ఉపయోగించాలి; ఈ సందర్భంలో “$\cap$”, “$\cup$”ల నిర్వచనాలు. అందువల్ల (a)ను ఇలా వాదిస్తాం:

\begin{quote}
(a) ముందుగా $X \cap (X \cup Y) \subseteq X$ అని చూపాలి. అంటే $\subseteq$ నిర్వచనం ప్రకారం, ఏ~$z$కైనా $z \in X \cap (X \cup Y)$ అయితే $z \in X$ అని చూపాలి. కాబట్టి $z \in X \cap (X \cup Y)$ అనుకుందాం. $z$ రెండు సమితుల ఖండనలో \tetoken{మూలకం}{!!{element}} అవడం అది రెండు సమితులలోనూ \tetoken{మూలకం}{!!{element}} అవడానికి అవసరమైనదీ చాలినదీ. అందువల్ల $z \in X$, అలాగే $z \in X \cup Y$ అని తేలుతుంది. ముఖ్యంగా మనం చూపాలనుకున్నట్లే $z \in X$.
\end{quote}

దీనితో నిరూపణ మొదటి సగం పూర్తయింది. చివరి దశలో ఒక పరికల్పన నుంచి ఒక సంయోగం ($z \in X$ మరియు $z \in X \cup Y$) వస్తే, అదే పరికల్పన నుంచి దాని ప్రతి సంయుక్త భాగం వస్తుంది అనే విషయాన్ని ఉపయోగించాం. ఈ నియమం మీకు “సంయోగ నిర్మూలనం” లేదా $\land$Elimగా తెలిసి ఉండవచ్చు. ఇప్పుడు (b)ను నిరూపిద్దాం:

\begin{quote}
(b) ఇప్పుడు $X \subseteq X \cap (X \cup Y)$ అని నిరూపిద్దాం. అంటే $\subseteq$ నిర్వచనం ప్రకారం, ఏ~$z$కైనా $z \in X$ అయితే $z \in X \cap (X \cup Y)$ అని చూపాలి. $z \in X$ అనుకుందాం. $z \in X \cap (X \cup Y)$ అని చూపడానికి, “$\cap$” నిర్వచనం ప్రకారం (i) $z \in X$, అలాగే (ii) $z \in X \cup Y$ అని చూపాలి. (i) మన పరికల్పనే, కాబట్టి ఇంకా నిరూపించాల్సిందేమీ లేదు. (ii) కోసం, $z$ సమితుల సమ్మేళనానికి \tetoken{మూలకం}{!!{element}} అవడం అది ఆ సమితుల్లో కనీసం ఒకదానికి \tetoken{మూలకం}{!!{element}} అవడానికి అవసరమైనదీ చాలినదీ అని గుర్తుంచుకోండి. $z \in X$; $X \cup Y$ అనేది $X$, $Y$ల సమ్మేళనం కాబట్టి ఇక్కడ ఆ నిబంధన తీరుతుంది. అందువల్ల $z \in X \cup Y$. (i) $z \in X$, (ii) $z \in X \cup Y$ రెండింటినీ చూపాం. కాబట్టి “$\cap$” నిర్వచనం ప్రకారం $z \in X \cap (X \cup Y)$.
\end{quote}

ఈ వివరణ కొంచెం విస్తారంగా ఉన్నా, సమితులు, వాటి సంబంధాల గురించి ఎలా తర్కిస్తామో చూపిస్తుంది. సాధారణంగా మనం ఇంత వివరంగా చెప్పం; ముఖ్యంగా అన్ని నిర్వచనాలను మళ్లీ మళ్లీ ప్రస్తావించకపోవచ్చు. మరింత ఉన్నత పాఠ్యంలో ఈ ఫలితం యొక్క నిరూపణ చాలా సంక్షిప్తంగా ఉంటుంది. ఉదాహరణకు ఇలా:
\end{explain}

\begin{prop}[శోషణ]
అన్ని సమితులు $X$, $Y$లకూ,
\[
X \cap (X \cup Y) = X
\]
\end{prop}

\begin{proof}
(a) $z \in X \cap (X \cup Y)$ అనుకుందాం. అప్పుడు $z \in X$; కాబట్టి $X \cap (X \cup Y) \subseteq X$.

(b) ఇప్పుడు $z \in X$ అనుకుందాం. అప్పుడు $z \in X \cup Y$ కూడా; అందువల్ల $z \in X \cap (X \cup Y)$ కూడా. కాబట్టి $X \subseteq X \cap (X \cup Y)$.
\end{proof}

\begin{prob}
$X \cup (X \cap Y) = X$ అని వివరంగా నిరూపించండి. తర్వాత సంక్షిప్త నిరూపణ ఇవ్వండి. (సూచన: $X \cup (X \cap Y) \subseteq X$ దిశలో సందర్భాలవారీ నిరూపణ, అంటే $\lor$Elim, అవసరం.)
\end{prob}

\end{document}
